\section{Introduction}

%Mapper is an efficient algorithm in the field of TDA, application + math developments.
%Its main weakness is its lack of robustness : many parameters too choose; cf papers.



\paragraph*{Mapper graphs and TDA.} The Mapper graph introduced in \cite{singh2007topological} is an essential tool of Topological Data Analysis (TDA), and has been used many times for visualization purposes
on different types of data, including, but not limited to, single-cell sequencing~\cite{wang2018topological,zechel2014topographical}, neural network architectures~\cite{9605336,joseph2021topological}, or 3D meshes~\cite{wang2020exploration,rosen2018inferring}.
Moreover, its ability to precisely encode (within the graph) the presence and sizes of geometric and topological structures in the data in a mathematically founded way
(through the use of algebraic topology) has also proved beneficial for highlighting subpopulations of interest, which are usually detected as peculiar
topological structures of significant sizes, and identifying the key features that best explain such subpopulations against the rest of the Mapper graph.
This general pipeline has become a key component in, e.g., biological inference in single-cell data sets, where differentiating stem cells can usually
be recovered from branching patterns in the corresponding Mapper graphs~\cite{Rizvi2017}.

\paragraph*{Parameter selection.} However, it has quickly become clear that the Mapper graph is quite sensitive to its parameters, in the
sense that the structure of the graph can vary a lot under (even small) changes of its parameters. As such, several pipelines based on Mapper
graphs actually involve brute force optimization: they first compute a grid of Mapper graphs corresponding to many different sets of parameters,
and then they pick the best one, either by manual inspection or with arbitrary criteria---leading to prohibitive running times.
In order to deal with this issue, several methods have been proposed in the literature for either assessing the statistical robustness of a given Mapper graph with respect to the distribution of the studied dataset~\cite{belchi2020numerical, Brown2021}, or for tuning the Mapper parameters automatically~\cite{carriere2018statistical}. 
Unfortunately, most tuning methods
involve simple heuristics that only work for some, but not all Mapper parameters; in particular, the so-called {\em filter parameter} has never been treated, to the best of our knowledge. This is mostly because it is a general continuous function, and can thus vary in a much wilder parameter space than the other Mapper parameters.

In another line of work, ensemble methods have recently been proposed to combine Mapper graphs over multiple parameter sets, rather than trying to find the best one \cite{kang2021ensemble,fitzpatrick2023ensemble},
%This empirically ensemble methods are known
so as to be able to produce outputs that are more robust. However, this imposes to aggregate families of completely different filter functions, with no guarantees on the resulting graph. %will lead  to a relevant final Mapper. 
In this work, we follow a different approach, and rather attempt at identifying an "optimal" filter function by minimizing specific loss functions.


Another approach related to our work is \cite{bui2020f}, where an alternative way of constructing Mapper graphs is proposed using a fuzzy clustering algorithm. %Our framework shares with F-Mapper 
Even though we also adopt a probabilistic approach (that allows, e.g., a point to belong to disconnected intervals in the cover of the filter range), the underlying probabilistic formalism 
that we use is new, while there is none in \cite{bui2020f}. In particular, we introduce stochastic {\it assignment schemes} and we address the parameter selection problem within this framework. 




 



\paragraph*{Contributions.} Our contribution is three-fold:
\begin{itemize}
    \item We introduce {\em Soft Mapper}: a generalization of the combinatorial Mapper graph in the form of a probability distribution on Mapper graphs,
% define a generalized Mapper as a distribution on Mapper (attention,  c'est bien defini ? cf mesurabilité dans papier avec Steve ???), which provides a smoothed  version of the original Mapper we call the Soft Mapper. 
    \item We propose a filter optimization framework adapted to a smooth Soft Mapper distribution with provable convergence guarantees,
%In this framework, we  propose an optimization method to find a filter function
    \item We implement and showcase the efficiency of Mapper filter optimization through Soft Mapper on various data sets, with public, open-source code in \texttt{TensorFlow}. 
\end{itemize}

The following of the article is organized as follows: in Section~\ref{sec:background} we recall the basics on the Mapper algorithm, then in Section~\ref{sec:soft_mapper} we detail the Soft Mapper construction, which is the main  focus of this work. We provide several interesting special cases of Soft Mapper in Section~\ref{sec:examples}, before introducing topological losses that are specific to Mapper graphs in Section~\ref{sec:tpological_loss}. We then present our optimization setting, in which a parameterized family of Mapper filter functions is optimized, in Section~\ref{sec:optim}, and we 
apply it on
%finally show the results of using our method in order to derive an optimal filters for 
3-dimensional shapes and single cell RNA-sequencing data in Section~\ref{sec:experiments}. Finally, we discuss the results of this article and present possible future work directions in Section~\ref{sec:discussion}.

